\scanpage{89}
De (vi) et (viii) résulte que pour tt motif $M$ et entier $n$, l'isomorphisme
\[
\unsure{\Phi|\mathscr{L}|\mathcal{L}}^s(M)\simeq\deleted{\unsure{M}}\Phi^s(M(n))
\]
est compatible avec les cônes amples.

De (iv) et (vi) résulte que si $p=0$, une forme ample sur $M$ induit une forme \uline{définie positive}
\[
\Gamma_k M\times\Gamma_k M\unsure{\mathbin{-}|\longrightarrow}\Gamma_k\mathbb{Q}(0)=\mathbb{Q}
\]

Si on décompose $M$ (de poids $p$) en composantes isotypiques $M_\alpha$, il vient
\[
\Phi^s(M)\simeq\mathop{\deleted{\unsure{\sum}}\unsure{\bigoplus|\prod}}\limits_\alpha\Phi^s(M_\alpha)
\]
donc pour étudier les formes positives, il suffit de les étudier en principe sur \unsure{des|les} morceaux isotypiques, qui sont d'ailleurs de la forme
\[
I\otimes_A\deleted{\unsure{M|N}}V
\]
où $I$ est un motif simple, $A$ son anneau des endomorphismes, qui est un corps, et $\deleted{M}V$ un \uline{vectoriel} sur $A$
% archive sheet number: 88
